\documentclass[12pt]{article}
\usepackage[utf8]{inputenc}
\usepackage{amsmath, amssymb, amsthm}
\usepackage{geometry}
\usepackage{fancyhdr}

\geometry{letterpaper, margin=1in}
\setlength{\headheight}{15pt}

% Header setup
\pagestyle{fancy}
\fancyhf{}
\lhead{Math 477 Spring 2026}
\rhead{Jeffrey Lu}
\cfoot{\thepage}

\title{Homework Assignment 1 \\ \large Math 292, Spring 2026 \\ Rutgers University}
\author{Jeffrey Lu}

\begin{document}

\maketitle

\noindent \textbf{Problem 1:}

\begin{enumerate}
      \item[(a)]
            \begin{align*}
                  \text{P(At least 1 pair)}1 - \frac{\binom{10}{6} * 2^6}{\binom{20}{6}} \approx{0.6533}
            \end{align*}
      \item[(b)]
            \begin{align*}
                  \text{P(Exactly 2 pair)} = \frac{\binom{10}{2}\binom{8}{2}2^2}{\binom{20}{6}}
            \end{align*}
            Then,
            \begin{align*}
                  \text{P(Exactly 2 pair $|$ At least 1 pair)} & = \frac{\text{P(Exactly 2 pair $\cap$ At least 1 pair)}}{\text{P(At least 1 pair)}} \\
                                                               & = \frac{5040}{25340}                                                                \\
                                                               & \approx{0.1991}
            \end{align*}

\end{enumerate}

\noindent \textbf{Problem 2:}
We can use total probability, either the second card is heart suited or it isn't
\begin{align*}
      P(H_1) & =  P(H_1~|~H_2 \cap{} H_3  \cap{} H_4) * P(H_2 | H_3 \cap{} H_4) + P(H_1~|~H_2^c \cap{} H_3  \cap{} H_4) * P(H_2^c | H_3 \cap{} H_4) \\
             & = \frac{10}{49}*\frac{11}{50} + \frac{11}{49}*\frac{39}{50}                                                                          \\
             & = 0.22
\end{align*}


\noindent \textbf{Problem 3:}

The only chance for HHHH to happen is at the very start since for it to happen after that, T must precede it. Therefore, with the law of total probability, we conclude
\begin{align*}
      \text{P(THHH before HHHH)} & = 1 - P(HHHH) \\ &= 1 - p^4
\end{align*}

\noindent \textbf{Problem 4:}



\begin{enumerate}
      \item[(a)] Let A = Suspect A is the burglar, B = Suspect B is the burglar, G = Glove is at scene
            \begin{align*}
                  P(A~|~G) & = \frac{P(G~|~A)P(A)}{P(G~|~A)P(A) + P(G~|~B)P(B)} \\
                           & = \frac{1 * 0.5}{1 * 0.5 + 0.5 * 0.1}              \\
                           & \approx{0.909}
            \end{align*}
      \item[(b)] Let O = Suspect B owns the glov
            \begin{align*}
                  P(O ~|~ G) & = \frac{P(G ~|~ O)P(O)}{P(G)} \\
                             & = \frac{1 * 0.1}{0.55}        \\
                             & \approx{0.182}
            \end{align*}

\end{enumerate}

\noindent \textbf{Problem 5:}
\begin{align*}
      P(x=0) & = P(N_1 < N_2) = 1/2                                   \\
      P(x=1) & = P(N_1 > N_2) * P(N_1 < N_3) = 1/2 * 1/3 = 1/6        \\
      P(x=2) & = P(N_1 > N_2, N_3) * P(N_1 < N_4) = 1/3 * 1/4 =  1/12 \\
      P(x=3) & = (3!)/5! = 1/20                                       \\
      P(x=4) & = P(N_1 = 5) =  1/5
\end{align*}



\end{document}